% !TEX program = xelatex
\documentclass[UTF8,12pt,a4paper]{ctexart}

% =========================================================
% 通用期末模拟卷 LaTeX 模板
% 用途：根据样卷题型和分值结构生成模拟卷及答案解析
% 敏感信息原则：模板内只保留占位符，不保留个人、学校或课程内部信息
% =========================================================

\usepackage{geometry}
\usepackage{array,booktabs,longtable,multirow,tabularx}
\usepackage{amsmath,amssymb,amsthm,bm,mathtools}
\usepackage{enumitem}
\usepackage{xcolor}
\usepackage{fancyhdr}
\usepackage{lastpage}
\usepackage{hyperref}
\usepackage{fvextra}
\usepackage{eso-pic}
\usepackage{tikz}

\geometry{left=2cm,right=2cm,top=2cm,bottom=1.8cm}

% =========================
% 字体设置
% =========================
\setmainfont{Noto Serif CJK SC}
\setCJKmainfont{Noto Serif CJK SC}
\setmonofont{DejaVu Sans Mono}

% =========================
% 可修改变量
% =========================
\newcommand{\CourseName}{课程名称}
\newcommand{\ExamSetName}{一/二}
\newcommand{\SchoolInfo}{学校/学院：未明确}
\newcommand{\TermInfo}{学年学期：未明确}
\newcommand{\MajorInfo}{专业/年级：未明确}
\newcommand{\ExamType}{闭卷/开卷：未明确}
\newcommand{\ExamDuration}{考试时间：未明确}
\newcommand{\ExamScore}{100 分}
\newcommand{\ExamLevel}{适用层级：未明确}
\newcommand{\WatermarkText}{仅供个人复习使用｜请勿商用传播}

% 是否启用水印：true/false
\newif\ifshowwatermark
\showwatermarktrue

% =========================
% 页眉页脚
% =========================
\pagestyle{fancy}
\fancyhf{}
\fancyhead[C]{\small 《\CourseName》期末模拟卷\ExamSetName}
\fancyfoot[C]{\small 第 \thepage\ 页 / 共 \pageref{LastPage} 页}
\renewcommand{\headrulewidth}{0pt}
\renewcommand{\footrulewidth}{0pt}

% =========================
% 水印
% =========================
\ifshowwatermark
\AddToShipoutPictureBG{%
  \begin{tikzpicture}[remember picture,overlay]
    \node[rotate=35,scale=1.6,text opacity=0.10] at (current page.center) {\WatermarkText};
  \end{tikzpicture}%
}
\fi

% =========================
% 列表与代码块
% =========================
\setlist[enumerate,1]{leftmargin=2em,itemsep=0.4em}
\setlist[enumerate,2]{leftmargin=2em,itemsep=0.25em}
\DefineVerbatimEnvironment{Code}{Verbatim}{fontsize=\small,breaklines,breakanywhere,frame=single,framesep=2mm}

% =========================
% 常用环境
% =========================
\newenvironment{choices}{\begin{enumerate}[label=(\Alph*),leftmargin=3.5em,itemsep=0.1em]}{\end{enumerate}}

\newcommand{\examtitle}{%
\begin{center}
  {\LARGE\bfseries 《\CourseName》本科期末模拟卷\ExamSetName}\\[0.4em]
  {\small \SchoolInfo\quad \TermInfo\quad \MajorInfo}\\[0.8em]
\end{center}}

\newcommand{\examinfo}{%
\begin{center}
\begin{tabular}{|>{\centering\arraybackslash}p{2.5cm}|>{\centering\arraybackslash}p{3.5cm}|>{\centering\arraybackslash}p{2.5cm}|>{\centering\arraybackslash}p{3.5cm}|}
\hline
考试方式 & \ExamType & 考试时间 & \ExamDuration \\ \hline
卷面总分 & \ExamScore & 适用层级 & \ExamLevel \\ \hline
\end{tabular}
\end{center}}

% 按样卷实际题型增删列，务必保证总分一致
\newcommand{\scoretable}{%
\begin{center}
\begin{tabular}{|c|c|c|c|c|c|}
\hline
题号 & 一 & 二 & 三 & 四 & 总分 \\ \hline
得分 & & & & & \\ \hline
\end{tabular}
\end{center}}

\begin{document}
\examtitle
\examinfo
\scoretable

\noindent\textbf{命题说明：}
本卷题型、题量和分值应严格参考样卷；题目必须重新命制，不能直接复制旧题、只改数字、只换变量名或只打乱选项。若样卷结构不同，请按样卷调整下面各大题。

% =========================================================
\section*{一、单项选择题（每小题 3 分，共 30 分）}
% =========================================================

\begin{enumerate}
\item 【题干】（\quad）
\begin{choices}
\item 【选项 A】
\item 【选项 B】
\item 【选项 C】
\item 【选项 D】
\end{choices}

\item 【题干】（\quad）
\begin{choices}
\item 【选项 A】
\item 【选项 B】
\item 【选项 C】
\item 【选项 D】
\end{choices}

\item 【题干】（\quad）
\begin{choices}
\item 【选项 A】
\item 【选项 B】
\item 【选项 C】
\item 【选项 D】
\end{choices}

% 继续补足至样卷规定题量
\end{enumerate}

% =========================================================
\section*{二、填空题（每小题 3 分，共 15 分）}
% =========================================================

\begin{enumerate}
\item 【题干】答案：\underline{\hspace{3cm}}。
\item 【题干】答案：\underline{\hspace{3cm}}。
\item 【题干】答案：\underline{\hspace{3cm}}。
% 继续补足至样卷规定题量
\end{enumerate}

% =========================================================
\section*{三、计算题 / 程序题 / 简答题（每小题按样卷分值设置，共若干分）}
% =========================================================

\begin{enumerate}
\item 【题干】
\begin{enumerate}[label=(\arabic*)]
\item 【小问】
\item 【小问】
\end{enumerate}

\item 【题干】
\begin{enumerate}[label=(\arabic*)]
\item 【小问】
\item 【小问】
\end{enumerate}

% 继续补足至样卷规定题量
\end{enumerate}

% =========================================================
\section*{四、综合题 / 应用题 / 论述题（共若干分）}
% =========================================================

\begin{enumerate}
\item 【题干】
\begin{enumerate}[label=(\arabic*)]
\item 【小问】
\item 【小问】
\end{enumerate}
\end{enumerate}

\newpage

% =========================================================
\section*{答案与解析}
% =========================================================

\subsection*{一、单项选择题}
\begin{enumerate}
\item 答案：【A/B/C/D】。解析：【写出判断依据；必要时说明排除理由】。
\item 答案：【A/B/C/D】。解析：【写出判断依据；必要时说明排除理由】。
\item 答案：【A/B/C/D】。解析：【写出判断依据；必要时说明排除理由】。
\end{enumerate}

\subsection*{二、填空题}
\begin{enumerate}
\item 答案：【结果】。解析：【关键计算、定义、规则或记忆点】。
\item 答案：【结果】。解析：【关键计算、定义、规则或记忆点】。
\item 答案：【结果】。解析：【关键计算、定义、规则或记忆点】。
\end{enumerate}

\subsection*{三、计算题 / 程序题 / 简答题}
\begin{enumerate}
\item
\textbf{解题思路：}【说明题型识别与方法选择】。\par
\textbf{关键公式/规则：}【列出公式、规则、代码逻辑或论述框架】。\par
\textbf{解答过程：}【完整推导、计算、代码输出或分点说明】。\par
\textbf{最终答案：}【结果】。\par
\textbf{易错提醒：}【提醒扣分点】。

\item
\textbf{解题思路：}【说明题型识别与方法选择】。\par
\textbf{关键公式/规则：}【列出公式、规则、代码逻辑或论述框架】。\par
\textbf{解答过程：}【完整推导、计算、代码输出或分点说明】。\par
\textbf{最终答案：}【结果】。\par
\textbf{易错提醒：}【提醒扣分点】。
\end{enumerate}

\subsection*{四、综合题 / 应用题 / 论述题}
\begin{enumerate}
\item
\textbf{解题思路：}【说明建模、分析或论证路线】。\par
\textbf{关键公式/规则：}【列出公式、规则或观点框架】。\par
\textbf{主要过程：}【完整过程】。\par
\textbf{最终答案：}【结果】。\par
\textbf{易错提醒：}【提醒扣分点】。
\end{enumerate}

\end{document}
